% VI24_FULL_SOLUTION_REFINEMENT_V1
\subsection*{VI.24.P4: Identity and iterated base change}
\begin{problem}\label{prob:vi24-p04}
Prove $X_S\cong X$ and $(X_{S'})_{S''}\cong X_{S''}$ for $S''\to S'\to S$.
\end{problem}

\ifdefined\FullProblemDossiers
\paragraph{Why this problem matters.}
This dossier isolates a reusable piece of the base-change calculus and records both its categorical and affine content.

\paragraph{Useful ideas before starting.}
Start from the Cartesian universal property.  When the schemes are affine, translate the square contravariantly into tensor products, quotients, or localizations, and then translate the resulting ring statement back into geometry.

\paragraph{Guided hints.}
\begin{enumerate}
\item Identify the relevant pullback or scalar-extension ring.
\item Use the appropriate universal property or exactness statement.
\item Check that the canonical maps are the expected projections.
\item State the geometric consequence, not only the ring isomorphism.
\end{enumerate}

\begin{solution}
For identity base change, a $T$-point of $X\times_SS$ is a pair $(a,b)$ with $fa=b$, so $b$ is determined by $a$.  Hence $X\times_SS$ represents the same functor as $X$, and the canonical isomorphism is induced by the first projection.

For iteration, expand the left side:
\[
(X_{S'})_{S''}=(X\times_SS')\times_{S'}S''.
\]
A map from $T$ to this scheme is a triple of maps to $X,S',S''$ in which the $S'$-component is forced to be the composite $T\to S''\to S'$, and the map to $X$ has the same composite to $S$.  Therefore the datum is exactly a compatible pair $T\to X$, $T\to S''$, which is represented by $X\times_SS''=X_{S''}$.  Naturality gives the canonical isomorphism.  Affine-locally this is the associativity isomorphism $(B\otimes_AA')\otimes_{A'}A''\cong B\otimes_AA''$.
\end{solution}

\paragraph{What happened mathematically.}
The base-changed object is forced by a universal property; tensor products make that universal property computable on affine charts.

\paragraph{Extensions.}
The same pattern will be reused in VI/25 for fibers and in VI/26 for properties of morphisms under change of base.

\paragraph{Provenance.}
Canonical functoriality dossier.
\else
\paragraph{Provenance.} Canonical functoriality dossier.
\fi
